\section{Conclusion: What Remains Open}

I began this paper with the sense that these names were doing more than producing isolated gematria values. After recalculation, source checking, expansion searches, and matched controls, I still think there is a pattern worth presenting. I would now describe it with more care and, in some places, with more wonder.

The strongest results are compact enough to see without specialized mathematics:

\begin{itemize}
\item YHWH becomes the base-12 palindrome $22_{12}$.
\item Shaddai becomes $222_{12}$ while retaining the decimal digits 314; its shared YHWH frame with Yeshua makes their exact decimal difference 72.
\item Ehyeh Asher Ehyeh becomes $393_{12}$, while its decimal digits sum to 12, its two outer names total 42, and its visible frame is YHWH Echad at $33_{12}$.
\item Elohim becomes $72_{12}$ and separately gives the decimal palindrome 646 under Mispar Gadol.
\item El Roi and El Gibbor share the palindromic value 242, which reaches 8.
\item El Olam crosses decimal and base-12 palindrome layers under its two value systems.
\item Yeshua becomes $282_{12}$: terminal 8, exact sum 12, and a palindrome that places 8 inside the visible YHWH frame, echoing John 14:11.
\item Yeshua HaMashiach becomes $525_{12}$, with the article completing both the palindrome and the exact sum of 12 inside the $55_{12}$ Adonai frame.
\item Across all 132 repeated-frame palindromes, $282_{12}$ and $525_{12}$ are the only two whose center equals the decimal digital root and whose digits sum to 12.
\end{itemize}

At the canonical level, 22, 42, and 72 recur through different kinds of contact rather than one proven sequence. Psalm 119 also joins 22 and 8 in a textual structure: twenty-two alphabetic units of eight verses each. That correspondence does not identify decimal 22 with $22_{12}$, but it gives the visible relationship another canonical setting. The chapter total 929 converts exactly to $655_{12}$, and one declared palindromic closure produces $6556_{12}$. Its notation collapses through decimal 22 to 4; its converted decimal value sums separately to 12. I no longer describe that object as though every step were one inevitable conversion. I still find the completed structure striking.

The audit also changed the paper in necessary ways. Some names had to be corrected. Some proposed titles were not titles in their cited verses. Some equalities disappeared when an article was restored. Base 3, not base 12, has the highest raw palindrome count in the repository-locked primary corpus. The matched controls do not make the primary base-12 count statistically decisive. The old language of a 12 attractor was too strong. Yeshua's value is not lexically unique, named outer frames occur in 23 of the 30 eligible value layers in the full expanded registry, the five anchors are not statistically privileged among all attainable alternatives, and the $X$--Asher--$X$ mechanism can produce other palindromes.

None of those corrections turns $22_{12}$ back into 26 on the page, removes the 8 from $282_{12}$, breaks $525_{12}$, creates a third solution in the complete 132-member family, or makes the complete Ehyeh formula behave like Ehyeh alone. Nor do the controls reproduce the identity and scriptural setting of the names whose arithmetic they imitate. They change the scale of the conclusions, not the existence of the calculations.

One robustness result deserves to be read from both directions. When Yeshua is removed, most of the network's cross-name connectivity disappears. As a statistical statement, the network depends heavily on one node. Within Christian theology, Yeshua standing at the connective center may be exactly what gives the architecture coherence. I do not think the numerical test can choose between those descriptions, and I do not want the paper to choose for the reader.

This is where I want the paper to stop telling the reader what to think. The numerical architecture may reflect features of Hebrew morphology, the way the corpus was selected, traditional development, deliberate arrangement, providential order, chance, or several of these at once. Mathematics can test some parts of that list more readily than others. It cannot collapse all of them into one final digit.

My own interest lies in the convergence. Yeshua places 8 inside the visible YHWH signature 22, visually echoes the reciprocal indwelling language of John 14:11, sums to 12, becomes the connective center of the tested network, and yields a matched collision set partitioned through 12, 8, and 4. Yeshua HaMashiach places its own decimal root inside the Adonai signature 55 and recalls the confession that Jesus Christ is Lord. Together they occupy the only two solutions in the complete repeated-frame family. Shaddai shares Yeshua's YHWH frame and stands exactly 72 away. A complete divine declaration acquires a form its repeated component does not possess and carries YHWH Echad as its frame. An added title moves 314 to 345, from a $\pi$ echo to a Pythagorean one. A canonical construction passes through 22 and ends at 4. Each observation is modest in isolation. Together, they form the architecture explored here.

I offer that architecture as a finding, not a verdict. The appendices explain every step, and the repository preserves the complete rows. What the pattern finally means is left where I believe it belongs: with the reader.
